
\documentclass{article}
\usepackage{colortbl}
\usepackage{makecell}
\usepackage{multirow}
\usepackage{supertabular}

\begin{document}

\newcounter{utterance}

\centering \large Interaction Transcript for game `taboo', experiment `medium\_en', episode 12 with reflexion{-}llama{-}3.1{-}8b{-}instruct.
\vspace{24pt}

{ \footnotesize  \setcounter{utterance}{1}
\setlength{\tabcolsep}{0pt}
\begin{supertabular}{c@{$\;$}|p{.15\linewidth}@{}p{.15\linewidth}p{.15\linewidth}p{.15\linewidth}p{.15\linewidth}p{.15\linewidth}}
    \# & $\;$A & \multicolumn{4}{c}{Game Master} & $\;\:$B\\
    \hline

    \theutterance \stepcounter{utterance}  
    & & \multicolumn{4}{p{0.6\linewidth}}{
        \cellcolor[rgb]{0.9,0.9,0.9}{
            \makecell[{{p{\linewidth}}}]{
                \texttt{\tiny{[P1$\langle$GM]}}
                \texttt{You have attempted this same task before and did not succeed. These are the plans you wrote after those attempts.} \\
\texttt{Trial \#1: New plan: I failed to prevent the other player from guessing the target word in three tries. Upon reflection, I realize that I should have provided more subtle and indirect clues that don't directly relate to the target word or its related words. In the first clue, I mentioned "population's characteristics," which is a direct reference to the target word. In the second clue, I mentioned "planning and policy{-}making," which is a common application of demographic data, but still too closely related. In the third clue, I mentioned "number of individuals in a specific area," which is a characteristic of demographic data, but again too direct.} \\
\\ 
\texttt{Next time, I will aim to provide more abstract and indirect clues that still convey the essence of the target word without giving away its identity. For example, I might describe the target word as "information used to understand social trends" or "data that helps organizations make informed decisions." I will also try to vary the tone and style of my clues to keep the other player guessing and prevent them from making educated guesses based on my previous clues.} \\
\\ 
\texttt{Trial \#2: New plan: I failed to prevent the other player from guessing the target word in three tries. Upon reflection, I realize that I should have provided more subtle and indirect clues that don't directly relate to the target word or its related words. In the first clue, I mentioned "population's characteristics," which is a direct reference to the target word. In the second clue, I mentioned "planning and policy{-}making," which is a common application of demographic data, but still too closely related. In the third clue, I mentioned "number of individuals in a specific area," which is a characteristic of demographic data, but again too direct.} \\
\\ 
\texttt{Next time, I will aim to provide more abstract and indirect clues that still convey the essence of the target word without giving away its identity. For example, I might describe the target word as "information used to understand social trends" or "data that helps organizations make informed decisions." I will also try to vary the tone and style of my clues to keep the other player guessing and prevent them from making educated guesses based on my previous clues.} \\
\\ 
\texttt{To achieve this, I will:} \\
\\ 
\texttt{1. Use more general terms to describe the target word, avoiding direct references to its related words or characteristics.} \\
\texttt{2. Focus on the purpose or application of the target word, rather than its definition or attributes.} \\
\texttt{3. Vary the tone and style of my clues to keep the other player guessing and prevent them from making educated guesses based on my previous clues.} \\
\texttt{4. Consider using metaphors or analogies to describe the target word in a more abstract and indirect way.} \\
\\ 
\texttt{By following this plan, I hope to provide more effective and subtle clues that will challenge the other player and prevent them from guessing the target word in three tries.} \\
\\ 
\\ 
\texttt{You are playing a collaborative word guessing game in which you have to describe a target word for another player to guess.} \\
\\ 
\texttt{Rules:} \\
\texttt{(a) You have to reply in the form: CLUE: <some text>. Guesses from the other player will start with GUESS.} \\
\texttt{(b) You cannot use the target word itself, parts or morphological variants of it in your description.} \\
\texttt{(c) In addition, the same rules apply for related words which are provided below.} \\
\\ 
\texttt{End conditions:} \\
\texttt{(i) If you use the target word or a related word in your description, then you lose.} \\
\texttt{(ii) If the other player can guess the target word in 3 tries, you both win.} \\
\\ 
\texttt{Let us start.} \\
\\ 
\texttt{This is the target word that you need to describe and that the other player needs to guess:} \\
\\ 
\texttt{demographic} \\
\\ 
\texttt{Related words are:} \\
\\ 
\texttt{{-} people} \\
\texttt{{-} group} \\
\texttt{{-} statistic} \\
\\ 
\texttt{Important: You are under time pressure, give short descriptions that are to the point!} \\
            }
        }
    }
    & \\ \\

    \theutterance \stepcounter{utterance}  
    & \multicolumn{4}{p{0.6\linewidth}}{
        \cellcolor[rgb]{0.9,0.9,0.9}{
            \makecell[{{p{\linewidth}}}]{
                \texttt{\tiny{[P1$\rangle$GM]}}
                \texttt{CLUE: Numbers that help us understand how a community is changing over time.} \\
            }
        }
    }
    & & \\ \\

    \theutterance \stepcounter{utterance}  
    & & & \multicolumn{2}{p{0.3\linewidth}}{
        \cellcolor[rgb]{0.9,0.9,0.9}{
            \makecell[{{p{\linewidth}}}]{
                \texttt{\tiny{[GM$|$GM]}}
                \texttt{continue} \\
            }
        }
    }
    & & \\ \\

    \theutterance \stepcounter{utterance}  
    & & & \multicolumn{2}{p{0.3\linewidth}}{
        \cellcolor[rgb]{0.9,0.9,0.9}{
            \makecell[{{p{\linewidth}}}]{
                \texttt{\tiny{[GM$|$GM]}}
                \texttt{numbers that help us understand how a community is changing over time} \\
            }
        }
    }
    & & \\ \\

    \theutterance \stepcounter{utterance}  
    & & \multicolumn{4}{p{0.6\linewidth}}{
        \cellcolor[rgb]{0.9,0.9,0.9}{
            \makecell[{{p{\linewidth}}}]{
                \texttt{\tiny{[GM$\rangle$P2]}}
                \texttt{You have attempted this same task before and did not succeed. These are the plans you wrote after those attempts.} \\
\texttt{Trial \#1: New plan: } \\
\\ 
\texttt{In my last attempt, I was able to make three correct guesses, but I failed to guess the target word. Upon reflection, I realize that I was focusing too much on the individual clues and not enough on the overall pattern of the game. The game seems to be building on the previous clues, with each new clue providing more specific information about the target word.} \\
\\ 
\texttt{Next time, I will pay closer attention to the sequence of clues and try to identify a pattern or a connection between them. I will also try to think more broadly about the possible meanings of the target word, rather than just focusing on the individual clues. For example, I might consider how the clues relate to each other, or how they might be connected to a broader concept or idea.} \\
\\ 
\texttt{In this case, the first clue mentioned "data about a population's characteristics," which led me to guess "STATISTICS." The second clue mentioned "used in planning and policy{-}making," which led me to guess "DEMOGRAPHY." The third clue mentioned "relates to the number of individuals in a specific area," which led me to guess "POPULATION."} \\
\\ 
\texttt{Looking back, I realize that I was trying to guess a word that was too specific, and I didn't consider the broader connections between the clues. Next time, I will try to think more holistically about the game and look for a word that encompasses the overall pattern of the clues.} \\
\\ 
\texttt{My next guess will be a word that takes into account the connections between the clues and tries to capture the broader meaning of the target word.} \\
\\ 
\texttt{Trial \#2: New plan: } \\
\\ 
\texttt{In my last attempt, I was able to make two correct guesses, but I failed to guess the target word. Upon reflection, I realize that I was not following the specified format for making a guess. The format is "GUESS: <a word>", but I was only providing "GUESS: <a word>" without the colon.} \\
\\ 
\texttt{Next time, I will make sure to follow the correct format for making a guess, including the colon before the word. For example, my next guess will be "GUESS: INFORMATION".} \\
\\ 
\\ 
\texttt{You are playing a collaborative word guessing game in which you have to guess a target word that another player describes to you.} \\
\\ 
\texttt{You can make one guess at each trial. You win when you guess the target word. You lose when you cannot guess it in 3 tries.} \\
\\ 
\texttt{After each trial you will get a new hint from the other player which starts with CLUE.} \\
\\ 
\texttt{Make your guesses by just saying the word using the following form: GUESS: <a word>} \\
\\ 
\texttt{Let us start.} \\
\\ 
\\ 
\texttt{CLUE: Numbers that help us understand how a community is changing over time.} \\
            }
        }
    }
    & \\ \\

    \theutterance \stepcounter{utterance}  
    & & & \multicolumn{4}{p{0.6\linewidth}}{
        \cellcolor[rgb]{0.9,0.9,0.9}{
            \makecell[{{p{\linewidth}}}]{
                \texttt{\tiny{[GM$\langle$P2]}}
                \texttt{GUESS: TRENDS} \\
            }
        }
    }
     \\ \\

    \theutterance \stepcounter{utterance}  
    & & & \multicolumn{2}{p{0.3\linewidth}}{
        \cellcolor[rgb]{0.9,0.9,0.9}{
            \makecell[{{p{\linewidth}}}]{
                \texttt{\tiny{[GM$|$GM]}}
                \texttt{continue} \\
            }
        }
    }
    & & \\ \\

    \theutterance \stepcounter{utterance}  
    & & & \multicolumn{2}{p{0.3\linewidth}}{
        \cellcolor[rgb]{0.9,0.9,0.9}{
            \makecell[{{p{\linewidth}}}]{
                \texttt{\tiny{[GM$|$GM]}}
                \texttt{trends} \\
            }
        }
    }
    & & \\ \\

    \theutterance \stepcounter{utterance}  
    & & \multicolumn{4}{p{0.6\linewidth}}{
        \cellcolor[rgb]{0.9,0.9,0.9}{
            \makecell[{{p{\linewidth}}}]{
                \texttt{\tiny{[P1$\langle$GM]}}
                \texttt{GUESS: TRENDS} \\
            }
        }
    }
    & \\ \\

    \theutterance \stepcounter{utterance}  
    & \multicolumn{4}{p{0.6\linewidth}}{
        \cellcolor[rgb]{0.9,0.9,0.9}{
            \makecell[{{p{\linewidth}}}]{
                \texttt{\tiny{[P1$\rangle$GM]}}
                \texttt{CLUE: A snapshot of a community's characteristics, often used to inform big decisions.} \\
            }
        }
    }
    & & \\ \\

    \theutterance \stepcounter{utterance}  
    & & & \multicolumn{2}{p{0.3\linewidth}}{
        \cellcolor[rgb]{0.9,0.9,0.9}{
            \makecell[{{p{\linewidth}}}]{
                \texttt{\tiny{[GM$|$GM]}}
                \texttt{continue} \\
            }
        }
    }
    & & \\ \\

    \theutterance \stepcounter{utterance}  
    & & & \multicolumn{2}{p{0.3\linewidth}}{
        \cellcolor[rgb]{0.9,0.9,0.9}{
            \makecell[{{p{\linewidth}}}]{
                \texttt{\tiny{[GM$|$GM]}}
                \texttt{a snapshot of a communitys characteristics often used to inform big decisions} \\
            }
        }
    }
    & & \\ \\

    \theutterance \stepcounter{utterance}  
    & & \multicolumn{4}{p{0.6\linewidth}}{
        \cellcolor[rgb]{0.9,0.9,0.9}{
            \makecell[{{p{\linewidth}}}]{
                \texttt{\tiny{[GM$\rangle$P2]}}
                \texttt{CLUE: A snapshot of a community's characteristics, often used to inform big decisions.} \\
            }
        }
    }
    & \\ \\

    \theutterance \stepcounter{utterance}  
    & & & \multicolumn{4}{p{0.6\linewidth}}{
        \cellcolor[rgb]{0.9,0.9,0.9}{
            \makecell[{{p{\linewidth}}}]{
                \texttt{\tiny{[GM$\langle$P2]}}
                \texttt{GUESS: PROFILE} \\
            }
        }
    }
     \\ \\

    \theutterance \stepcounter{utterance}  
    & & & \multicolumn{2}{p{0.3\linewidth}}{
        \cellcolor[rgb]{0.9,0.9,0.9}{
            \makecell[{{p{\linewidth}}}]{
                \texttt{\tiny{[GM$|$GM]}}
                \texttt{continue} \\
            }
        }
    }
    & & \\ \\

    \theutterance \stepcounter{utterance}  
    & & & \multicolumn{2}{p{0.3\linewidth}}{
        \cellcolor[rgb]{0.9,0.9,0.9}{
            \makecell[{{p{\linewidth}}}]{
                \texttt{\tiny{[GM$|$GM]}}
                \texttt{profile} \\
            }
        }
    }
    & & \\ \\

    \theutterance \stepcounter{utterance}  
    & & \multicolumn{4}{p{0.6\linewidth}}{
        \cellcolor[rgb]{0.9,0.9,0.9}{
            \makecell[{{p{\linewidth}}}]{
                \texttt{\tiny{[P1$\langle$GM]}}
                \texttt{GUESS: PROFILE} \\
            }
        }
    }
    & \\ \\

    \theutterance \stepcounter{utterance}  
    & \multicolumn{4}{p{0.6\linewidth}}{
        \cellcolor[rgb]{0.9,0.9,0.9}{
            \makecell[{{p{\linewidth}}}]{
                \texttt{\tiny{[P1$\rangle$GM]}}
                \texttt{CLUE: A tool for city planners, helping them make sense of a city's growth and needs.} \\
            }
        }
    }
    & & \\ \\

    \theutterance \stepcounter{utterance}  
    & & & \multicolumn{2}{p{0.3\linewidth}}{
        \cellcolor[rgb]{0.9,0.9,0.9}{
            \makecell[{{p{\linewidth}}}]{
                \texttt{\tiny{[GM$|$GM]}}
                \texttt{continue} \\
            }
        }
    }
    & & \\ \\

    \theutterance \stepcounter{utterance}  
    & & & \multicolumn{2}{p{0.3\linewidth}}{
        \cellcolor[rgb]{0.9,0.9,0.9}{
            \makecell[{{p{\linewidth}}}]{
                \texttt{\tiny{[GM$|$GM]}}
                \texttt{a tool for city planners helping them make sense of a citys growth and needs} \\
            }
        }
    }
    & & \\ \\

    \theutterance \stepcounter{utterance}  
    & & \multicolumn{4}{p{0.6\linewidth}}{
        \cellcolor[rgb]{0.9,0.9,0.9}{
            \makecell[{{p{\linewidth}}}]{
                \texttt{\tiny{[GM$\rangle$P2]}}
                \texttt{CLUE: A tool for city planners, helping them make sense of a city's growth and needs.} \\
            }
        }
    }
    & \\ \\

    \theutterance \stepcounter{utterance}  
    & & & \multicolumn{4}{p{0.6\linewidth}}{
        \cellcolor[rgb]{0.9,0.9,0.9}{
            \makecell[{{p{\linewidth}}}]{
                \texttt{\tiny{[GM$\langle$P2]}}
                \texttt{GUESS: ANALYSIS} \\
            }
        }
    }
     \\ \\

    \theutterance \stepcounter{utterance}  
    & & & \multicolumn{2}{p{0.3\linewidth}}{
        \cellcolor[rgb]{0.9,0.9,0.9}{
            \makecell[{{p{\linewidth}}}]{
                \texttt{\tiny{[GM$|$GM]}}
                \texttt{continue} \\
            }
        }
    }
    & & \\ \\

    \theutterance \stepcounter{utterance}  
    & & & \multicolumn{2}{p{0.3\linewidth}}{
        \cellcolor[rgb]{0.9,0.9,0.9}{
            \makecell[{{p{\linewidth}}}]{
                \texttt{\tiny{[GM$|$GM]}}
                \texttt{analysis} \\
            }
        }
    }
    & & \\ \\

    \theutterance \stepcounter{utterance}  
    & & & \multicolumn{2}{p{0.3\linewidth}}{
        \cellcolor[rgb]{0.9,0.9,0.9}{
            \makecell[{{p{\linewidth}}}]{
                \texttt{\tiny{[GM$|$GM]}}
                \texttt{3} \\
            }
        }
    }
    & & \\ \\

    \theutterance \stepcounter{utterance}  
    & \multicolumn{4}{p{0.6\linewidth}}{
        \cellcolor[rgb]{0.95,0.95,0.95}{
            \makecell[{{p{\linewidth}}}]{
                \texttt{\tiny{[P2$\rightarrow$P1]}}
                \texttt{Post{-}episode Reflexion (not sent to GM):} \\
\texttt{New plan: I failed to prevent the other player from guessing the target word in three tries. Upon reflection, I realize that I should have provided more subtle and indirect clues that don't directly relate to the target word or its related words. In the first clue, I mentioned "population's characteristics," which is a direct reference to the target word. In the second clue, I mentioned "planning and policy{-}making," which is a common application of demographic data, but still too closely related. In the third clue, I mentioned "number of individuals in a specific area," which is a characteristic of demographic data, but again too direct.} \\
\\ 
\texttt{Next time, I will aim to provide more abstract and indirect clues that still convey the essence of the target word without giving away its identity. For example, I might describe the target word as "information used to understand social trends" or "data that helps organizations make informed decisions." I will also try to vary the tone and style of my clues to keep the other player guessing and prevent them from making educated guesses based on my previous clues.} \\
\\ 
\texttt{To achieve this, I will:} \\
\\ 
\texttt{1. Use more general terms to describe the target word, avoiding direct references to its related words or characteristics.} \\
\texttt{2. Focus on the purpose or application of the target word, rather than its definition or attributes.} \\
\texttt{3. Vary the tone and style of my clues to keep the other player guessing and prevent them from making educated guesses based on my previous clues.} \\
\texttt{4. Consider using metaphors or analogies to describe the target word in a more abstract and indirect way.} \\
\\ 
\texttt{Additionally, I will pay closer attention to the related words provided and try to avoid using language that is too closely associated with them. I will also try to create a sense of mystery and intrigue around the target word, making it more challenging for the other player to guess.} \\
\\ 
\texttt{For example, my next clue might be: "A type of information that helps us understand how a community is changing, but is often hidden in plain sight." This clue still conveys the essence of the target word, but in a more abstract and indirect way, and avoids direct references to its related words or characteristics.} \\
            }
        }
    }
    & & \\ \\

    \theutterance \stepcounter{utterance}  
    & & & \multicolumn{4}{p{0.6\linewidth}}{
        \cellcolor[rgb]{0.95,0.95,0.95}{
            \makecell[{{p{\linewidth}}}]{
                \texttt{\tiny{[P1$\rightarrow$P2]}}
                \texttt{Post{-}episode Reflexion (not sent to GM):} \\
\texttt{New plan: } \\
\\ 
\texttt{In my last attempt, I was able to make three correct guesses, but I failed to guess the target word. Upon reflection, I realize that I was focusing too much on the individual clues and not enough on the overall pattern of the game. The game seems to be building on the previous clues, with each new clue providing more specific information about the target word.} \\
\\ 
\texttt{However, I also realize that I was trying to guess a word that was too specific, and I didn't consider the broader connections between the clues. Next time, I will try to think more holistically about the game and look for a word that encompasses the overall pattern of the clues.} \\
\\ 
\texttt{In addition, I will pay closer attention to the sequence of clues and try to identify a pattern or a connection between them. I will also try to think more broadly about the possible meanings of the target word, rather than just focusing on the individual clues.} \\
\\ 
\texttt{My next guess will be a word that takes into account the connections between the clues and tries to capture the broader meaning of the target word. I will also make sure to follow the correct format for making a guess, including the colon before the word.} \\
\\ 
\texttt{For example, my next guess will be "GUESS: DEMOGRAPHICS".} \\
            }
        }
    }
     \\ \\

\end{supertabular}
}

\end{document}
